% !TEX root = ../高等概率论.tex
% =====================================================================
%  随机级数：0-1 律，Kolmogorov 极大不等式，随机级数的收敛，三级数定理，Kronecker 引理，收敛速度，无穷期望
% =====================================================================
\lecture{尾 \texorpdfstring{$\sigma$}{σ}-代数与 Kolmogorov \texorpdfstring{$0$-$1$}{0-1} 律，Hewitt--Savage \texorpdfstring{$0$-$1$}{0-1} 律，Kolmogorov 极大不等式，独立随机变量级数的收敛，三级数定理，Kronecker 引理，强大数定律的第二个证明，收敛速度，期望为无穷的情形}
\label{lec:random-series}

考虑随机符号的调和级数
\[
\sum_{n=1}^\infty\frac{\varepsilon_n}n,
\]
其中 $\varepsilon_1,\varepsilon_2,\dots$ 独立，各以概率 $\frac12$ 取 $\pm1$。若符号全取 $+1$，这是发散的调和级数；若正负交替，这是收敛到 $\log2$ 的交错级数。随机地选符号，级数收敛的概率是多少？这一讲的回答是：概率是 $1$。而且这个问题的答案只能是 $0$ 或 $1$，不可能是 $\frac12$ 之类的数，因为“级数是否收敛”不依赖于有限个项。

这一讲研究独立随机变量之和 $\sum_nX_n$ 何时几乎必然收敛，并由此给出强大数定律的第二个证明。这条路线比上一讲 Etemadi 的证明长，但它有两个好处：它给出收敛的速度（比如 $S_n/n$ 以多快的速度趋于 $\mu$），并且能处理期望为无穷的情形。教材把这一节标为选读；本讲中标有 ★ 的部分第一遍可以跳过，跳过它们不影响后面的课程。

\topic{本讲用到的前文结果}

\begin{enumerate}
\item 独立性的 $\pi$-类判别（定理~\ref{thm:independence:pi-systems}）与分组定理（定理~\ref{thm:independence:grouping}）；若事件 $A$ 与自己独立，则 $P(A)\in\{0,1\}$（\egref{ex:independence:constant}）。
\item 独立随机变量乘积的期望等于期望的乘积（定理~\ref{thm:independence:expectation}）；独立随机变量和的方差等于方差的和。
\item 两个 Borel--Cantelli 引理（第~\ref{lec:borel-cantelli} 讲）；强大数定律需要有限的期望（定理~\ref{thm:borel-cantelli:necessity}）。
\item 上一讲的截断引理~\ref{lem:strong-laws:truncation} 与方差估计（引理~\ref{lem:strong-laws:variance}）：对 i.i.d.\ 且 $E\abs{X_1}<\infty$，$Y_k=X_k\one_{\{\abs{X_k}\leqslant k\}}$ 满足 $\sum_k\Var(Y_k)/k^2\leqslant4E\abs{X_1}$。
\item 尾概率公式（定理~\ref{thm:product:tail-formula}）与 Tonelli 定理。
\item 生成的 $\sigma$-代数、好集原理与 $\pi$-$\lambda$ 定理（第~\ref{lec:set-systems} 讲）；测度的唯一性定理（定理~\ref{thm:measures:uniqueness-finite}）。
\end{enumerate}

\topic{尾 \texorpdfstring{$\sigma$}{σ}-代数与 Kolmogorov \texorpdfstring{$0$-$1$}{0-1} 律}

设 $X_1,X_2,\dots$ 是随机变量。令
\[
\cF_n'=\sigma(X_n,X_{n+1},\dots),\qquad\cT=\bigcap_{n\geqslant1}\cF_n' .
\]
$\cF_n'$ 是“第 $n$ 时刻以后的未来”，$\cT$ 称为\term{尾 $\sigma$-代数}（tail $\sigma$-field），是“遥远的未来”。直观上，$A\in\cT$ 当且仅当改变有限个 $X_i$ 的值不影响 $A$ 是否发生。

\begin{example}\label{ex:random-series:tail-io}
若 $B_n\in\cR$，则 $\{X_n\in B_n\ \text{i.o.}\}\in\cT$。事实上对每个 $m$，$\{X_n\in B_n\ \text{i.o.}\}=\bigcap_{k\geqslant m}\bigcup_{n\geqslant k}\{X_n\in B_n\}\in\cF_m'$（上极限可以从任何位置开始算）。取 $X_n=\one_{A_n}$，$B_n=\{1\}$，就得到 $\{A_n\ \text{i.o.}\}\in\cT$（这里 $\cT$ 是 $\one_{A_n}$ 的尾 $\sigma$-代数）。
\end{example}

\begin{example}\label{ex:random-series:tail-sums}
令 $S_n=X_1+\cdots+X_n$。
\begin{enumerate}
\item $\{\lim_nS_n\ \text{存在}\}\in\cT$：对每个 $m$，$\lim_nS_n$ 存在当且仅当 $\lim_n(S_n-S_m)$ 存在，而 $S_n-S_m=X_{m+1}+\cdots+X_n$ 是 $\cF_{m+1}'$-可测的。
\item 若 $c_n\to\infty$，则 $\{\limsup_nS_n/c_n>x\}\in\cT$：对固定的 $m$，$S_m/c_n\to0$，所以 $\limsup_nS_n/c_n=\limsup_n(S_n-S_m)/c_n$，后者是 $\cF_{m+1}'$-可测的。
\item $\{\limsup_nS_n>0\}$ 一般不属于 $\cT$。例如 $X_1$ 以概率 $\frac12$ 取 $\pm1$，$X_n=0$（$n\geqslant2$）。则 $\cF_n'=\{\varnothing,\Omega\}$（$n\geqslant2$），$\cT=\{\varnothing,\Omega\}$，而 $\{\limsup S_n>0\}=\{X_1=1\}$ 的概率是 $\frac12$，不属于 $\cT$。改变 $X_1$ 会改变这个事件是否发生。
\end{enumerate}
\end{example}

下面的定理说，对独立的随机变量，尾 $\sigma$-代数中的事件都是平凡的。

\begin{theorem}[Kolmogorov $0$-$1$ 律]\label{thm:random-series:kolmogorov}
设 $X_1,X_2,\dots$ 独立，$A\in\cT$。则 $P(A)=0$ 或 $1$。
\end{theorem}

\begin{remark}[证明思路]
我们证明 $A$ 与自己独立：$P(A)=P(A\cap A)=P(A)^2$，从而 $P(A)\in\{0,1\}$。分两步“包抄”：(a) 过去 $\sigma(X_1,\dots,X_k)$ 与未来 $\sigma(X_{k+1},\dots)$ 独立；(b) 于是 $\cT$（它包含于每个未来之中）与每个过去独立，从而与一切过去的并独立，再由 $\pi$-$\lambda$ 定理，与 $\sigma(X_1,X_2,\dots)$ 独立。而 $\cT$ 本身包含于 $\sigma(X_1,X_2,\dots)$。
\end{remark}

\begin{proof}
\emph{(a) $\sigma(X_1,\dots,X_k)$ 与 $\sigma(X_{k+1},X_{k+2},\dots)$ 独立。}由分组定理（定理~\ref{thm:independence:grouping}），把 $\{\sigma(X_i)\}_{i\geqslant1}$ 分成 $\{1,\dots,k\}$ 与 $\{k+1,k+2,\dots\}$ 两组即得。（分组定理允许无穷多个指标。）

\emph{(b) $\sigma(X_1,X_2,\dots)$ 与 $\cT$ 独立。}令 $\cA=\bigcup_k\sigma(X_1,\dots,X_k)$。它是 $\pi$-类：$\sigma(X_1,\dots,X_k)$ 随 $k$ 递增，两个元素总属于同一个 $\sigma(X_1,\dots,X_k)$，交仍在其中。$\cT$ 是 $\sigma$-代数，当然是 $\pi$-类。若 $B\in\sigma(X_1,\dots,X_k)$，$C\in\cT\subseteq\sigma(X_{k+1},\dots)$，由 (a)，$P(B\cap C)=P(B)P(C)$。所以 $\cA$ 与 $\cT$ 独立。由定理~\ref{thm:independence:pi-systems}，$\sigma(\cA)$ 与 $\sigma(\cT)=\cT$ 独立。而 $\sigma(\cA)=\sigma(X_1,X_2,\dots)$：每个 $X_i$ 关于 $\sigma(\cA)$ 可测，并且 $\cA\subseteq\sigma(X_1,X_2,\dots)$。

\emph{(c)}$\cT\subseteq\cF_1'=\sigma(X_1,X_2,\dots)$，所以 $A\in\cT$ 也属于 $\sigma(X_1,X_2,\dots)$，由 (b)，$A$ 与 $A$ 独立，$P(A)=P(A)^2$。
\end{proof}

\begin{corollary}\label{cor:random-series:tail-constant}
设 $X_1,X_2,\dots$ 独立，$Y$ 是 $\cT$-可测的广义实值随机变量。则存在常数 $c\in[-\infty,\infty]$ 使 $Y=c$ a.s.。特别地，若 $c_n\to\infty$，则 $\limsup_nS_n/c_n$ 几乎必然等于一个常数；$P(\sum_nX_n\ \text{收敛})$ 只能是 $0$ 或 $1$。
\end{corollary}
\begin{proof}
对每个 $x$，$\{Y\leqslant x\}\in\cT$，概率是 $0$ 或 $1$。$F(x)=P(Y\leqslant x)$ 单调不减，只取 $0,1$ 两个值。令 $c=\sup\{x:F(x)=0\}$（若集合为空，$c=-\infty$；若为全体实数，$c=+\infty$）。$x<c$ 时 $P(Y\leqslant x)=0$，$x>c$ 时 $P(Y\leqslant x)=1$；取 $x=c\pm\frac1k$ 并令 $k\to\infty$，得 $P(Y<c)=0$，$P(Y\leqslant c)=1$（$c$ 有限时），即 $Y=c$ a.s.；$c=\pm\infty$ 的情形同理。其余两句由\egref{ex:random-series:tail-sums} 得到。
\end{proof}

例如，对独立的事件，$P(A_n\ \text{i.o.})\in\{0,1\}$，这与两个 Borel--Cantelli 引理一致；$0$-$1$ 律告诉我们答案只能是 $0$ 或 $1$，Borel--Cantelli 引理告诉我们是哪一个。

\topic{Hewitt--Savage \texorpdfstring{$0$-$1$}{0-1} 律★}

对 i.i.d.\ 序列，有一个更强的 $0$-$1$ 律。它允许事件依赖于每一个 $X_i$，只要求它在有限个 $X_i$ 的重排下不变。比如 $\{S_n\in B\ \text{i.o.}\}$：改变 $X_1$ 的值会改变它，所以它一般不在 $\cT$ 中；可是交换 $X_1$ 与 $X_2$ 不改变 $n\geqslant2$ 的 $S_n$，所以它在重排下不变。

为了说清楚“重排”，我们在序列空间上讨论。令 $\R^\N$ 为一切实数列 $x=(x_1,x_2,\dots)$ 的集合，$\cR^\N$ 为由坐标映射 $x\mapsto x_i$ 生成的 $\sigma$-代数（与第~\ref{lec:set-systems} 讲作业 D 中的 $\sigma(\cC)$ 相同）。$\N$ 的一个\term{有限置换}是一个双射 $\pi:\N\to\N$，只有有限个 $i$ 使 $\pi(i)\neq i$。对 $x\in\R^\N$，令 $(\pi x)_i=x_{\pi(i)}$。若 $B\in\cR^\N$ 对一切有限置换 $\pi$ 满足 $\{x:\pi x\in B\}=B$，称 $B$ 是\term{可交换的}（exchangeable）。

\begin{theorem}[Hewitt--Savage $0$-$1$ 律]\label{thm:random-series:hewitt-savage}
设 $X_1,X_2,\dots$ 是 i.i.d.，$\mathbf X=(X_1,X_2,\dots)$，$B\in\cR^\N$ 是可交换的。则 $P(\mathbf X\in B)\in\{0,1\}$。
\end{theorem}
\begin{proof}[证明★]
记 $A=\{\mathbf X\in B\}$，$\mu=P\circ\mathbf X^{-1}$ 为 $\mathbf X$ 在 $(\R^\N,\cR^\N)$ 上的分布（$\mathbf X$ 可测，因为每个坐标 $X_i$ 可测，第~\ref{lec:random-variables} 讲定理~\ref{thm:random-variables:generators}）。

\emph{第一步：用有限维事件逼近。}令 $\cC=\bigcup_n\{\{x:(x_1,\dots,x_n)\in D\}:D\in\cR^n\}$ 为柱集的全体，它是代数，生成 $\cR^\N$。令 $\cG$ 为这样的 $E\in\cR^\N$ 的全体：对任意 $\varepsilon>0$ 存在 $C\in\cC$ 使 $\mu(E\triangle C)<\varepsilon$。$\cC\subseteq\cG$。$\cG$ 对补集封闭（$E^c\triangle C^c=E\triangle C$）。对可数并：设 $E=\bigcup_iE_i$，$E_i\in\cG$。取 $N$ 使 $\mu(E\setminus\bigcup_{i\leqslant N}E_i)<\frac\varepsilon2$（从下连续），再取 $C_i\in\cC$ 使 $\mu(E_i\triangle C_i)<\frac\varepsilon{2N}$；$C=\bigcup_{i\leqslant N}C_i\in\cC$，并且 $E\triangle C\subseteq(E\setminus\bigcup_{i\leqslant N}E_i)\cup\bigcup_{i\leqslant N}(E_i\triangle C_i)$，$\mu(E\triangle C)<\varepsilon$。所以 $\cG$ 是包含 $\cC$ 的 $\sigma$-代数，$\cG=\cR^\N$。于是存在 $B_n\in\cR^n$，使 $A_n=\{(X_1,\dots,X_n)\in B_n\}$ 满足 $P(A_n\triangle A)\to0$。

\emph{第二步：重排后的分布不变。}固定 $n$，令 $\pi$ 交换 $j$ 与 $j+n$（$1\leqslant j\leqslant n$），其余不动，$\pi$ 是有限置换，$\pi\circ\pi$ 是恒等映射。随机序列 $\pi\mathbf X=(X_{\pi(1)},X_{\pi(2)},\dots)$ 的分量仍然 i.i.d.，所以对任意 $k$ 与 Borel 集 $D_1,\dots,D_k$，$P(X_{\pi(1)}\in D_1,\dots,X_{\pi(k)}\in D_k)=\prod_iP(X_1\in D_i)=P(X_1\in D_1,\dots,X_k\in D_k)$。形如 $\{x:x_1\in D_1,\dots,x_k\in D_k\}$ 的集合组成生成 $\cR^\N$ 的 $\pi$-类，由唯一性定理，$\pi\mathbf X$ 与 $\mathbf X$ 有相同的分布 $\mu$。

\emph{第三步：交换前后两段。}令 $E_n=\{x:(x_1,\dots,x_n)\in B_n\}\triangle B$。由第二步，$P(\pi\mathbf X\in E_n)=P(\mathbf X\in E_n)=P(A_n\triangle A)$。另一方面，$\{\pi\mathbf X\in B\}=\{\mathbf X\in B\}=A$（$B$ 可交换），而 $\{((\pi\mathbf X)_1,\dots,(\pi\mathbf X)_n)\in B_n\}=\{(X_{n+1},\dots,X_{2n})\in B_n\}=:A_n'$。所以 $\{\pi\mathbf X\in E_n\}=A_n'\triangle A$，于是 $P(A_n'\triangle A)=P(A_n\triangle A)\to0$。

\emph{第四步。}由 $\abs{P(E)-P(F)}\leqslant P(E\triangle F)$，$P(A_n)\to P(A)$，$P(A_n')\to P(A)$。由 $A_n\triangle A_n'\subseteq(A_n\triangle A)\cup(A\triangle A_n')$，$P(A_n\triangle A_n')\to0$，于是
\[
0\leqslant P(A_n)-P(A_n\cap A_n')\leqslant P(A_n\cup A_n')-P(A_n\cap A_n')=P(A_n\triangle A_n')\to0,
\]
$P(A_n\cap A_n')\to P(A)$。但 $A_n\in\sigma(X_1,\dots,X_n)$ 与 $A_n'\in\sigma(X_{n+1},\dots,X_{2n})$ 独立，$P(A_n\cap A_n')=P(A_n)P(A_n')\to P(A)^2$。所以 $P(A)=P(A)^2$。
\end{proof}

尾事件都是可交换的（重排有限个坐标，只影响有限个 $X_i$），所以对 i.i.d.\ 序列，Hewitt--Savage $0$-$1$ 律比 Kolmogorov $0$-$1$ 律强。比如对 i.i.d.\ 序列，$P(S_n\in B\ \text{i.o.})\in\{0,1\}$；以后讨论随机游走的常返性时要用到这一点。

\topic{Kolmogorov 极大不等式}

由 $0$-$1$ 律，$P(\sum_nX_n\ \text{收敛})$ 是 $0$ 或 $1$。下面的工具帮助我们证明它是 $1$。

\begin{theorem}[Kolmogorov 极大不等式]\label{thm:random-series:maximal}
设 $X_1,\dots,X_n$ 独立，$EX_i=0$，$\Var(X_i)<\infty$，$S_k=X_1+\cdots+X_k$。则对 $x>0$，
\[
P\Bigl(\max_{1\leqslant k\leqslant n}\abs{S_k}\geqslant x\Bigr)\leqslant\frac{\Var(S_n)}{x^2}.
\]
\end{theorem}

Chebyshev 不等式只给出 $P(\abs{S_n}\geqslant x)\leqslant\Var(S_n)/x^2$。极大不等式说，同样的界对整条路径的最大值也成立：路径中途“冲高”的概率，并不比终点“冲高”的概率大多少。

\begin{proof}
\emph{按第一次越过 $x$ 的时刻分解。}令 $A_k=\{\abs{S_k}\geqslant x,\ \abs{S_j}<x\ \text{对一切}\ j<k\}$，$k=1,\dots,n$。它们两两不交，并为 $\{\max_k\abs{S_k}\geqslant x\}$。

\emph{展开。}由 $(S_n-S_k)^2\geqslant0$，
\[
\begin{aligned}
ES_n^2&\geqslant\sum_{k=1}^nE(S_n^2;A_k)=\sum_{k=1}^nE\bigl(S_k^2+2S_k(S_n-S_k)+(S_n-S_k)^2;A_k\bigr)\\
&\geqslant\sum_{k=1}^nE(S_k^2;A_k)+\sum_{k=1}^nE\bigl(2S_k\one_{A_k}(S_n-S_k)\bigr).
\end{aligned}
\]
\emph{交叉项为零。}$S_k\one_{A_k}$ 是 $X_1,\dots,X_k$ 的函数，$S_n-S_k$ 是 $X_{k+1},\dots,X_n$ 的函数，由推论~\ref{cor:independence:functions} 它们独立，都平方可积从而可积，由定理~\ref{thm:independence:expectation}，$E(2S_k\one_{A_k}(S_n-S_k))=2E(S_k\one_{A_k})E(S_n-S_k)=0$。

\emph{结论。}在 $A_k$ 上 $\abs{S_k}\geqslant x$，所以
\[
\Var(S_n)=ES_n^2\geqslant\sum_{k=1}^nE(S_k^2;A_k)\geqslant x^2\sum_{k=1}^nP(A_k)=x^2P\Bigl(\max_{k\leqslant n}\abs{S_k}\geqslant x\Bigr).\qedhere
\]
\end{proof}

这个“按第一次越过的时刻分解”的技巧，以后在鞅论中会变成 Doob 极大不等式，那时 $A_k$ 是一个“停时”的事件。

\topic{独立随机变量级数的收敛}

\begin{theorem}\label{thm:random-series:convergence}
设 $X_1,X_2,\dots$ 独立，$EX_n=0$。若 $\sum_n\Var(X_n)<\infty$，则 $\sum_nX_n$ 几乎必然收敛。
\end{theorem}

\begin{remark}[证明思路]
级数收敛等价于部分和是 Cauchy 列，即“$M$ 以后的部分和的最大摆动”$w_M=\sup_{m,n\geqslant M}\abs{S_m-S_n}$ 趋于零。极大不等式正好控制最大摆动：$P(\sup_{m\geqslant M}\abs{S_m-S_M}>\varepsilon)\leqslant\varepsilon^{-2}\sum_{n>M}\Var(X_n)$，右边是收敛级数的余项。
\end{remark}

\begin{proof}
令 $S_N=\sum_{n\leqslant N}X_n$。对 $M<N$，把极大不等式用于 $X_{M+1},\dots,X_N$：
\[
P\Bigl(\max_{M\leqslant m\leqslant N}\abs{S_m-S_M}>\varepsilon\Bigr)\leqslant\frac{\Var(S_N-S_M)}{\varepsilon^2}=\frac1{\varepsilon^2}\sum_{n=M+1}^N\Var(X_n).
\]
令 $N\to\infty$，左边的事件递增趋于 $\{\sup_{m\geqslant M}\abs{S_m-S_M}>\varepsilon\}$，由从下连续，
\[
P\Bigl(\sup_{m\geqslant M}\abs{S_m-S_M}>\varepsilon\Bigr)\leqslant\frac1{\varepsilon^2}\sum_{n>M}\Var(X_n)\to0\quad(M\to\infty).
\]
令 $w_M=\sup_{m,n\geqslant M}\abs{S_m-S_n}$。由三角不等式 $w_M\leqslant2\sup_{m\geqslant M}\abs{S_m-S_M}$，所以 $P(w_M>2\varepsilon)\to0$。$w_M$ 随 $M$ 单调不增，极限 $w=\lim_Mw_M$ 存在，并且 $P(w>2\varepsilon)\leqslant P(w_M>2\varepsilon)\to0$，于是 $P(w>2\varepsilon)=0$。$\varepsilon$ 取遍 $\frac1j$，$w=0$ a.s.。而 $w(\omega)=0$ 说明 $S_n(\omega)$ 是 Cauchy 列，从而收敛。
\end{proof}

\begin{example}[随机符号的级数]\label{ex:random-series:random-signs}
设 $\varepsilon_n$ 独立，各以概率 $\frac12$ 取 $\pm1$，$\alpha>0$，$X_n=\varepsilon_nn^{-\alpha}$。$EX_n=0$，$\Var(X_n)=n^{-2\alpha}$。若 $\alpha>\frac12$，由定理~\ref{thm:random-series:convergence}，$\sum_n\varepsilon_nn^{-\alpha}$ 几乎必然收敛。特别地（$\alpha=1$），本讲开头的随机调和级数几乎必然收敛。注意绝对收敛 $\sum_n\abs{X_n}<\infty$ 当且仅当 $\alpha>1$：当 $\frac12<\alpha\leqslant1$ 时，级数收敛完全靠随机符号的正负相消。三级数定理（下面）的必要性部分将说明 $\alpha>\frac12$ 也是必要的：$\alpha\leqslant\frac12$ 时级数几乎必然发散。
\end{example}

\begin{workedexample}\label{we:random-series:sparse-signs}
设 $X_n$ 独立，$P(X_n=1)=P(X_n=-1)=\frac{p_n}2$，$P(X_n=0)=1-p_n$。证明 $\sum_nX_n$ 几乎必然收敛当且仅当 $\sum_np_n<\infty$，并且 $\sum_np_n=\infty$ 时它几乎必然发散。
\end{workedexample}
\begin{solution}
若 $\sum_np_n<\infty$，则 $\sum_nP(X_n\neq0)<\infty$，由第一 Borel--Cantelli 引理只有有限个 $X_n\neq0$，级数几乎必然收敛（甚至只有有限个非零项）。也可以用定理~\ref{thm:random-series:convergence}：$EX_n=0$，$\Var(X_n)=p_n$。

若 $\sum_np_n=\infty$，事件 $\{X_n\neq0\}$ 独立，由第二 Borel--Cantelli 引理，几乎必然有无穷多个 $n$ 使 $\abs{X_n}=1$；级数的项不趋于零，级数发散。

与\egref{ex:random-series:random-signs} 比较：那里每一项都非零，但大小趋于零；这里每一项的大小要么是 $0$ 要么是 $1$，收敛只能靠“非零项只有有限个”。
\end{solution}

下面的定理给出了独立随机变量级数收敛的充要条件。

\begin{theorem}[Kolmogorov 三级数定理]\label{thm:random-series:three-series}
设 $X_1,X_2,\dots$ 独立，$A>0$，$Y_n=X_n\one_{\{\abs{X_n}\leqslant A\}}$。$\sum_nX_n$ 几乎必然收敛的充要条件是下面三个级数都收敛：
\[
\text{(i)}\ \sum_nP(\abs{X_n}>A),\qquad\text{(ii)}\ \sum_nEY_n,\qquad\text{(iii)}\ \sum_n\Var(Y_n).
\]
\end{theorem}
\begin{proof}
我们只证充分性；必要性要用中心极限定理，留到那里再证。令 $\mu_n=EY_n$。$Y_n-\mu_n$ 独立，期望为零，由 (iii) 与定理~\ref{thm:random-series:convergence}，$\sum_n(Y_n-\mu_n)$ 几乎必然收敛；再由 (ii)，$\sum_nY_n$ 几乎必然收敛。由 (i) 与第一 Borel--Cantelli 引理，$P(X_n\neq Y_n\ \text{i.o.})=0$，几乎必然地两个级数只有有限项不同，所以 $\sum_nX_n$ 几乎必然收敛。
\end{proof}

三个条件的含义：(i) 大的项只有有限个；(ii) 截断后的项“平均”起来收敛；(iii) 截断后的项的随机波动可以求和。

\topic{Kronecker 引理与强大数定律的第二个证明}

随机级数的收敛怎样联系到大数定律？下面这个关于数列的引理是桥梁。

\begin{theorem}[Kronecker 引理]\label{thm:random-series:kronecker}
设 $a_n>0$，$a_n\uparrow\infty$，$\sum_nx_n/a_n$ 收敛。则 $\frac1{a_n}\sum_{m=1}^nx_m\to0$。
\end{theorem}
\begin{proof}
令 $a_0=0$，$b_0=0$，$b_m=\sum_{k\leqslant m}x_k/a_k$，则 $x_m=a_m(b_m-b_{m-1})$。分部求和（Abel 变换）：
\[
\begin{aligned}
\frac1{a_n}\sum_{m=1}^nx_m&=\frac1{a_n}\Bigl(\sum_{m=1}^na_mb_m-\sum_{m=1}^na_mb_{m-1}\Bigr)=\frac1{a_n}\Bigl(a_nb_n+\sum_{m=2}^na_{m-1}b_{m-1}-\sum_{m=1}^na_mb_{m-1}\Bigr)\\
&=b_n-\sum_{m=1}^n\frac{a_m-a_{m-1}}{a_n}b_{m-1}.
\end{aligned}
\]
由假设 $b_n\to b_\infty$。最后的和是 $b_0,\dots,b_{n-1}$ 的加权平均（权重 $\frac{a_m-a_{m-1}}{a_n}\geqslant0$，总和为 $1$），直观上它也趋于 $b_\infty$。严格地：令 $B=\sup_n\abs{b_n}$，给定 $\varepsilon>0$，取 $M$ 使 $m\geqslant M$ 时 $\abs{b_m-b_\infty}<\frac\varepsilon2$，再取 $N$ 使 $n\geqslant N$ 时 $\frac{a_M}{a_n}<\frac\varepsilon{4B+1}$。则 $n\geqslant N$ 时
\[
\Bigl|\sum_{m=1}^n\frac{a_m-a_{m-1}}{a_n}b_{m-1}-b_\infty\Bigr|\leqslant\sum_{m=1}^n\frac{a_m-a_{m-1}}{a_n}\abs{b_{m-1}-b_\infty}\leqslant\frac{a_M}{a_n}\cdot2B+\frac{a_n-a_M}{a_n}\cdot\frac\varepsilon2<\varepsilon .
\]
（$m\leqslant M$ 的项权重之和是 $a_M/a_n$，每项 $\abs{b_{m-1}-b_\infty}\leqslant2B$；$m>M$ 的项 $\abs{b_{m-1}-b_\infty}<\frac\varepsilon2$。）所以 $\frac1{a_n}\sum_{m\leqslant n}x_m\to b_\infty-b_\infty=0$。
\end{proof}

\begin{theorem}[强大数定律，第二个证明]\label{thm:random-series:slln}
设 $X_1,X_2,\dots$ 是 i.i.d.，$E\abs{X_1}<\infty$，$\mu=EX_1$。则 $S_n/n\to\mu$ a.s.
\end{theorem}
\begin{proof}
令 $Y_k=X_k\one_{\{\abs{X_k}\leqslant k\}}$，$T_n=Y_1+\cdots+Y_n$。由引理~\ref{lem:strong-laws:truncation}，只需证 $T_n/n\to\mu$ a.s.。令 $Z_k=Y_k-EY_k$，$EZ_k=0$，$Z_k$ 独立。由引理~\ref{lem:strong-laws:variance}，
\[
\sum_k\Var\Bigl(\frac{Z_k}k\Bigr)=\sum_k\frac{\Var(Y_k)}{k^2}\leqslant4E\abs{X_1}<\infty .
\]
由定理~\ref{thm:random-series:convergence}，$\sum_kZ_k/k$ 几乎必然收敛；由 Kronecker 引理（$a_k=k$），$\frac1n\sum_{k\leqslant n}(Y_k-EY_k)\to0$ a.s.。由控制收敛定理 $EY_k\to\mu$，于是 Cesàro 平均 $\frac1n\sum_{k\leqslant n}EY_k\to\mu$，所以 $T_n/n\to\mu$ a.s.
\end{proof}

这个证明不需要把 $X_k$ 分成正部与负部，也不需要单调性和子列，代价是要求完全独立（极大不等式用到了独立性），而 Etemadi 的证明只要两两独立。

\topic{收敛的速度}

随机级数方法的一个好处是能给出收敛的速度。由于可以从每个 $X_i$ 中减去 $\mu$，下面不妨设 $\mu=0$。

\begin{theorem}\label{thm:random-series:rate-L2}
设 $X_1,X_2,\dots$ 是 i.i.d.，$EX_1=0$，$EX_1^2=\sigma^2<\infty$。则对任意 $\varepsilon>0$，
\[
\frac{S_n}{n^{1/2}(\log n)^{1/2+\varepsilon}}\to0\quad\text{a.s.}
\]
\end{theorem}
\begin{proof}
令 $a_n=n^{1/2}(\log n)^{1/2+\varepsilon}$（$n\geqslant2$），$a_1\in(0,a_2)$，则 $a_n$ 递增趋于无穷（$n\geqslant2$ 时 $n^{1/2}$ 与 $\log n$ 都递增）。
\[
\sum_{n\geqslant1}\Var\Bigl(\frac{X_n}{a_n}\Bigr)=\sigma^2\Bigl(\frac1{a_1^2}+\sum_{n\geqslant2}\frac1{n(\log n)^{1+2\varepsilon}}\Bigr)<\infty
\]
（与 $\int_2^\infty\frac{\dd x}{x(\log x)^{1+2\varepsilon}}<\infty$ 比较）。由定理~\ref{thm:random-series:convergence}，$\sum_nX_n/a_n$ 几乎必然收敛；由 Kronecker 引理即得。
\end{proof}

以后的重对数律会说明 $\limsup_nS_n/\sqrt{2n\log\log n}=\sigma$ a.s.，所以这个结论离最好的结果已经不远了：分母中的 $(\log n)^{1/2+\varepsilon}$ 只比最优的 $(\log\log n)^{1/2}$ 大一点。

\begin{theorem}[Marcinkiewicz--Zygmund]\label{thm:random-series:mz}
设 $X_1,X_2,\dots$ 是 i.i.d.，$EX_1=0$，$E\abs{X_1}^p<\infty$，$1<p<2$。则 $S_n/n^{1/p}\to0$ a.s.
\end{theorem}
\begin{proof}[证明★]
\emph{截断。}令 $Y_k=X_k\one_{\{\abs{X_k}\leqslant k^{1/p}\}}$，$T_n=Y_1+\cdots+Y_n$。$\sum_kP(Y_k\neq X_k)=\sum_kP(\abs{X_1}^p>k)\leqslant E\abs{X_1}^p<\infty$（命题~\ref{prop:expectation:tail-sum}），由第一 Borel--Cantelli 引理，只需证 $T_n/n^{1/p}\to0$。

\emph{方差。}与引理~\ref{lem:strong-laws:variance} 相同，
\[
\begin{aligned}
\sum_{m\geqslant1}\frac{\Var(Y_m)}{m^{2/p}}&\leqslant\sum_{m\geqslant1}\frac1{m^{2/p}}\int_0^{m^{1/p}}2yP(\abs{X_1}>y)\dd y\\
&=\sum_{n\geqslant1}\int_{(n-1)^{1/p}}^{n^{1/p}}2yP(\abs{X_1}>y)\Bigl(\sum_{m\geqslant n}m^{-2/p}\Bigr)\dd y,
\end{aligned}
\]
这里把 $[0,m^{1/p}]$ 分成 $[(n-1)^{1/p},n^{1/p}]$（$n\leqslant m$），再交换求和次序。因为 $\frac2p>1$，$\sum_{m\geqslant n}m^{-2/p}\leqslant n^{-2/p}+\int_n^\infty x^{-2/p}\dd x\leqslant Cn^{(p-2)/p}$；对 $y\in[(n-1)^{1/p},n^{1/p}]$，$n\geqslant y^p$，而 $(p-2)/p<0$，所以 $n^{(p-2)/p}\leqslant y^{p-2}$。于是上式不超过
\[
C\int_0^\infty2y\cdot y^{p-2}P(\abs{X_1}>y)\dd y=\frac{2C}p\int_0^\infty py^{p-1}P(\abs{X_1}>y)\dd y=\frac{2C}pE\abs{X_1}^p<\infty .
\]
由定理~\ref{thm:random-series:convergence} 与 Kronecker 引理（$a_m=m^{1/p}$），$n^{-1/p}\sum_{m\leqslant n}(Y_m-\mu_m)\to0$ a.s.，$\mu_m=EY_m$。

\emph{期望。}由 $EX_1=0$，$\mu_m=-E(X_1;\abs{X_1}>m^{1/p})$。在 $\{\abs{X_1}>m^{1/p}\}$ 上 $\frac{\abs{X_1}}{m^{1/p}}\leqslant\bigl(\frac{\abs{X_1}}{m^{1/p}}\bigr)^p$，所以
\[
\abs{\mu_m}\leqslant m^{1/p}E\Bigl(\frac{\abs{X_1}^p}m;\abs{X_1}>m^{1/p}\Bigr)=m^{-1+1/p}\delta_m,
\]
其中 $\delta_m=E(\abs{X_1}^p;\abs{X_1}^p>m)\to0$（控制收敛定理）。$\sum_{m\leqslant n}m^{-1+1/p}\leqslant pn^{1/p}$，与 Cesàro 平均的论证相同，$n^{-1/p}\sum_{m\leqslant n}m^{-1+1/p}\delta_m\to0$：给定 $\eta$，取 $M$ 使 $m>M$ 时 $\delta_m<\eta$，则上式 $\leqslant n^{-1/p}\sum_{m\leqslant M}m^{-1+1/p}\delta_m+p\eta$。所以 $T_n/n^{1/p}\to0$ a.s.
\end{proof}

$p=1$ 时这就是强大数定律（$\mu=0$）；$p$ 越接近 $2$，矩条件越强，速度 $n^{1/p}$ 越接近 $\sqrt n$。

\topic{期望为无穷的情形★}

St.\ Petersburg 游戏中 $EX_i=\infty$，$S_n/(n\log_2n)\toP1$，而 $\limsup_nS_n/(n\log_2n)=\infty$ a.s.（第~\ref{lec:borel-cantelli} 讲习题~\exref{ex:borel-cantelli:A13}）。下面 Feller 的定理说明，$E\abs{X_1}=\infty$ 时，$S_n/a_n$ 不可能几乎必然收敛到一个非零的极限。定理~\ref{thm:borel-cantelli:necessity} 是 $a_n=n$ 的特例。

\begin{theorem}[Feller]\label{thm:random-series:feller}
设 $X_1,X_2,\dots$ 是 i.i.d.，$E\abs{X_1}=\infty$，$a_n>0$ 且 $a_n/n$ 单调不减。则 $\limsup_n\abs{S_n}/a_n$ 几乎必然等于 $0$ 或 $\infty$，按 $\sum_nP(\abs{X_1}\geqslant a_n)$ 收敛或发散而定。
\end{theorem}
\begin{proof}[证明★]
$a_n/n$ 单调不减蕴含 $a_n$ 单调不减（$a_{n+1}\geqslant\frac{n+1}na_n$），并且对正整数 $k$，$a_{kn}\geqslant ka_n$。

\emph{发散的情形。}对正整数 $k$，由 $a_{kn}\geqslant ka_n$ 与 $a$ 单调，
\[
\sum_nP(\abs{X_1}\geqslant ka_n)\geqslant\sum_nP(\abs{X_1}\geqslant a_{kn})\geqslant\frac1k\sum_{m\geqslant k}P(\abs{X_1}\geqslant a_m)=\infty,
\]
（把 $m\in[kn,kn+k)$ 这 $k$ 项归到 $n$，每项 $\leqslant P(\abs{X_1}\geqslant a_{kn})$）。由第二 Borel--Cantelli 引理，$\abs{X_n}\geqslant ka_n$ 无穷多次，所以 $\limsup\abs{X_n}/a_n\geqslant k$，对一切 $k$ 成立。又 $\abs{X_n}\leqslant\abs{S_n}+\abs{S_{n-1}}\leqslant2\max\{\abs{S_n},\abs{S_{n-1}}\}$，$a_{n-1}\leqslant a_n$，所以 $\limsup\abs{S_n}/a_n=\infty$。

\emph{收敛的情形。}令 $a_0=0$。

(1) \emph{一个恒等式。}$a_n\geqslant na_1\to\infty$。把 $m=\sum_{n=1}^m1$ 代入并交换求和次序，
\begin{equation}\label{eq:random-series:feller-identity}
\sum_{m\geqslant1}mP(a_{m-1}\leqslant\abs{X_1}<a_m)=\sum_{n\geqslant1}P(\abs{X_1}\geqslant a_{n-1})<\infty,
\end{equation}
因为对固定的 $n$，区间 $[a_{m-1},a_m)$（$m\geqslant n$）不交，并为 $[a_{n-1},\infty)$；右边有限是假设。

(2) \emph{截断。}令 $Y_n=X_n\one_{\{\abs{X_n}<a_n\}}$，$T_n=Y_1+\cdots+Y_n$。由假设与第一 Borel--Cantelli 引理，只需证 $T_n/a_n\to0$。

(3) \emph{方差。}$\Var(Y_n/a_n)\leqslant a_n^{-2}\sum_{m=1}^nE(X_1^2;a_{m-1}\leqslant\abs{X_1}<a_m)$。交换求和次序，并用 $a_n\geqslant\frac nma_m$（$n\geqslant m$）得 $\sum_{n\geqslant m}a_n^{-2}\leqslant\frac{m^2}{a_m^2}\sum_{n\geqslant m}\frac1{n^2}\leqslant\frac{2m}{a_m^2}$，
\[
\sum_n\Var\Bigl(\frac{Y_n}{a_n}\Bigr)\leqslant\sum_mE(X_1^2;a_{m-1}\leqslant\abs{X_1}<a_m)\cdot\frac{2m}{a_m^2}\leqslant2\sum_mmP(a_{m-1}\leqslant\abs{X_1}<a_m)<\infty .
\]
由定理~\ref{thm:random-series:convergence} 与 Kronecker 引理，$(T_n-ET_n)/a_n\to0$ a.s.

(4) \emph{期望。}先说明 $a_n/n\to\infty$：否则 $a_n\leqslant Ln$，$\sum_nP(\abs{X_1}\geqslant Ln)=\infty$（因为 $E\abs{X_1}/L=\infty$），与假设矛盾。固定 $N$，对 $n>N$，
\[
\frac{\abs{ET_n}}{a_n}\leqslant\frac n{a_n}E(\abs{X_1};\abs{X_1}<a_n)\leqslant\frac{na_N}{a_n}+\frac n{a_n}E(\abs{X_1};a_N\leqslant\abs{X_1}<a_n).
\]
第一项 $\to0$。第二项：$m\leqslant n$ 时 $\frac n{a_n}\leqslant\frac m{a_m}$，在 $\{\abs{X_1}<a_m\}$ 上 $\frac m{a_m}\abs{X_1}\leqslant m$，所以第二项不超过
\[
\sum_{m=N+1}^n\frac m{a_m}E(\abs{X_1};a_{m-1}\leqslant\abs{X_1}<a_m)\leqslant\sum_{m>N}mP(a_{m-1}\leqslant\abs{X_1}<a_m),
\]
由 \eqref{eq:random-series:feller-identity}，它在 $N$ 充分大时任意小。所以 $ET_n/a_n\to0$，$T_n/a_n\to0$ a.s.
\end{proof}

\topic{本讲小结}

\begin{enumerate}
\item 尾 $\sigma$-代数 $\cT=\bigcap_n\sigma(X_n,X_{n+1},\dots)$ 中的事件不依赖于有限个变量。Kolmogorov $0$-$1$ 律：独立时尾事件的概率是 $0$ 或 $1$；尾可测的随机变量几乎必然是常数。对 i.i.d.\ 序列，可交换的事件也是如此（Hewitt--Savage）。
\item Kolmogorov 极大不等式：$P(\max_{k\leqslant n}\abs{S_k}\geqslant x)\leqslant\Var(S_n)/x^2$。证明的方法是按第一次越过 $x$ 的时刻分解。
\item 独立、期望为零、方差可求和，则 $\sum X_n$ 几乎必然收敛；三级数定理给出充要条件。随机调和级数 $\sum\pm\frac1n$ 几乎必然收敛。
\item Kronecker 引理把级数的收敛转化为平均的收敛：$\sum x_n/a_n$ 收敛蕴含 $\frac1{a_n}\sum_{m\leqslant n}x_m\to0$。由此得到强大数定律的第二个证明，以及收敛速度：方差有限时 $S_n=o(\sqrt n(\log n)^{1/2+\varepsilon})$，$E\abs X^p<\infty$ 时 $S_n=o(n^{1/p})$。
\item 期望为无穷时，$S_n/a_n$ 的上极限只能是 $0$ 或 $\infty$（Feller）。
\end{enumerate}

% =====================================================================
\homework{级数收敛的若干判别法，加权强大数定律，非负独立随机变量之和}[级数收敛的判别法，加权强大数定律，非负随机变量之和]
\label{hw:random-series}

\exercisepart{基本习题}

\begin{exercises}{A}{课堂内容的补充}
\item\label{ex:random-series:A1} 设 $X_1,X_2,\dots$ 是 i.i.d.，$EX_i=0$，$\Var(X_i)=C<\infty$，$p>\frac12$。不用随机级数，用极大不等式（取 $n_m=\lceil m^\alpha\rceil$，$\alpha(2p-1)>1$）证明 $S_n/n^p\to0$ a.s.
\begin{solution}
思路：在子列 $n_m$ 上用极大不等式，它同时控制了 $n_m$ 与 $n_{m+1}$ 之间的一切 $n$，所以不再需要单调性。

令 $n_m=\lceil m^\alpha\rceil$。由极大不等式，
\[
P\Bigl(\max_{k\leqslant n_{m+1}}\abs{S_k}\geqslant\varepsilon n_m^p\Bigr)\leqslant\frac{Cn_{m+1}}{\varepsilon^2n_m^{2p}}\leqslant\frac{C'}{\varepsilon^2}m^{\alpha(1-2p)}
\]
（$n_{m+1}\leqslant2(m+1)^\alpha\leqslant2^{1+\alpha}m^\alpha$，$n_m\geqslant m^\alpha$）。$\alpha(2p-1)>1$，可求和。由第一 Borel--Cantelli 引理，几乎必然地 $m$ 充分大时 $\max_{k\leqslant n_{m+1}}\abs{S_k}<\varepsilon n_m^p$。对 $n_m\leqslant n<n_{m+1}$，$\abs{S_n}/n^p\leqslant\max_{k\leqslant n_{m+1}}\abs{S_k}/n_m^p<\varepsilon$。$\varepsilon$ 取遍 $\frac1j$。
\end{solution}

\item\label{ex:random-series:A2} 设 $X_1,X_2,\dots$ 是 i.i.d.，$p>0$。证明：若 $S_n/n^{1/p}\to0$ a.s.，则 $E\abs{X_1}^p<\infty$。（这是定理~\ref{thm:random-series:mz} 的逆，容易得多。）
\begin{solution}
$\frac{X_n}{n^{1/p}}=\frac{S_n}{n^{1/p}}-\Bigl(\frac{n-1}n\Bigr)^{1/p}\frac{S_{n-1}}{(n-1)^{1/p}}\to0$ a.s.，所以 $P(\abs{X_n}>n^{1/p}\ \text{i.o.})=0$。事件 $\{\abs{X_n}^p>n\}$ 独立，由第二 Borel--Cantelli 引理，$\sum_nP(\abs{X_1}^p>n)<\infty$，由命题~\ref{prop:expectation:tail-sum}，$E\abs{X_1}^p<\infty$。
\end{solution}

\item\label{ex:random-series:A3} 设 $X_1,X_2,\dots$ 是 i.i.d.\ 的标准正态随机变量。证明对任意 $t\in\R$，$\sum_{n\geqslant1}X_n\frac{\sin(n\pi t)}n$ 几乎必然收敛。
\begin{solution}
各项独立，期望为零，方差 $\frac{\sin^2(n\pi t)}{n^2}\leqslant\frac1{n^2}$ 可求和，由定理~\ref{thm:random-series:convergence} 即得。（这个随机 Fourier 级数在适当的意义下给出 Brown 运动的一种构造，它的“几乎必然”是对每个固定的 $t$ 说的。）
\end{solution}

\item\label{ex:random-series:A4} 设 $X_n$ 独立，$EX_n=0$，$\Var(X_n)=\sigma_n^2$。
\begin{questions}
\item 证明：若 $\sum_n\sigma_n^2/n^2<\infty$，则 $\sum_nX_n/n$ 几乎必然收敛，从而 $\frac1n\sum_{m\leqslant n}X_m\to0$ a.s.
\begin{solution}
$X_n/n$ 独立，期望为零，方差 $\sigma_n^2/n^2$ 可求和，由定理~\ref{thm:random-series:convergence}，$\sum_nX_n/n$ 几乎必然收敛；由 Kronecker 引理（$a_n=n$）即得。这是“方差不同”的强大数定律，条件比 $L^2$ 弱大数定律的“方差有界”弱（只要 $\sigma_n^2=o(n/\log^{1+\varepsilon}n)$ 之类）。
\end{solution}
\item 设 $\sum_n\sigma_n^2/n^2=\infty$，不妨设 $\sigma_n^2\leqslant n^2$。证明存在独立的 $X_n$，$EX_n=0$，$\Var(X_n)\leqslant\sigma_n^2$，使 $X_n/n$ 从而 $\frac1n\sum_{m\leqslant n}X_m$ 不几乎必然收敛到 $0$。
\begin{solution}
令 $X_n$ 以概率 $\frac{\sigma_n^2}{2n^2}$ 各取 $\pm n$，以概率 $1-\frac{\sigma_n^2}{n^2}$ 取 $0$（$\sigma_n^2\leqslant n^2$ 保证这是概率分布）；由定理~\ref{thm:independence:existence}，这样的独立序列存在。$EX_n=0$，$\Var(X_n)=n^2\cdot\frac{\sigma_n^2}{n^2}=\sigma_n^2$。$\sum_nP(\abs{X_n}=n)=\sum_n\sigma_n^2/n^2=\infty$，由第二 Borel--Cantelli 引理，$\abs{X_n}/n=1$ 无穷多次，$X_n/n\not\to0$。又 $\frac{X_n}n=\frac{S_n}n-\frac{n-1}n\cdot\frac{S_{n-1}}{n-1}$，所以 $S_n/n$ 也不趋于零。这说明 A4-1) 的条件在“只知道方差”的情况下不能再减弱。
\end{solution}
\end{questions}

\item\label{ex:random-series:A5} 设 $\psi(x)=x^2$（$\abs x\leqslant1$），$\psi(x)=\abs x$（$\abs x\geqslant1$）。证明：若 $X_n$ 独立，$EX_n=0$，$\sum_nE\psi(X_n)<\infty$，则 $\sum_nX_n$ 几乎必然收敛。
\begin{solution}
用三级数定理（充分性），$A=1$，$Y_n=X_n\one_{\{\abs{X_n}\leqslant1\}}$。
\begin{itemize}
\item (i) $\abs x>1$ 时 $\psi(x)\geqslant1$，所以 $P(\abs{X_n}>1)\leqslant E\psi(X_n)$，可求和。
\item (iii) $\Var(Y_n)\leqslant EY_n^2=E(X_n^2;\abs{X_n}\leqslant1)\leqslant E\psi(X_n)$，可求和。
\item (ii) 由 $EX_n=0$，$EY_n=-E(X_n;\abs{X_n}>1)$，$\abs{EY_n}\leqslant E(\abs{X_n};\abs{X_n}>1)\leqslant E\psi(X_n)$，所以 $\sum_nEY_n$ 绝对收敛。
\end{itemize}
$\psi$ 在 $0$ 附近像 $x^2$（控制方差），在远处像 $\abs x$（控制大的值）。这道题统一了“方差可求和”与“$\sum E\abs{X_n}<\infty$”两种情形。
\end{solution}

\item\label{ex:random-series:A6} 设 $X_n$ 独立，$0<p(n)\leqslant2$，$p(n)>1$ 时 $EX_n=0$，并且 $\sum_nE\abs{X_n}^{p(n)}<\infty$。证明 $\sum_nX_n$ 几乎必然收敛。
\begin{solution}
三级数定理，$A=1$。(i) $P(\abs{X_n}>1)\leqslant E\abs{X_n}^{p(n)}$（Markov 不等式）。(iii) $\abs x\leqslant1$、$p\leqslant2$ 时 $x^2\leqslant\abs x^p$，所以 $\Var(Y_n)\leqslant EY_n^2\leqslant E\abs{X_n}^{p(n)}$。(ii) 若 $p(n)\leqslant1$：$\abs x\leqslant1$ 时 $\abs x\leqslant\abs x^{p(n)}$，$\abs{EY_n}\leqslant E\abs{X_n}^{p(n)}$。若 $p(n)>1$：$EY_n=-E(X_n;\abs{X_n}>1)$，而 $\abs x\geqslant1$ 时 $\abs x\leqslant\abs x^{p(n)}$，$\abs{EY_n}\leqslant E\abs{X_n}^{p(n)}$。所以三个级数都收敛。
\end{solution}

\item\label{ex:random-series:A7} （加权强大数定律）设 $X_1,X_2,\dots$ 是 i.i.d.，$E\abs{X_1}<\infty$，$\mu=EX_1$，$\{c_k\}$ 是有界的实数列，$\abs{c_k}\leqslant C$。证明 $\frac1n\sum_{k\leqslant n}c_k(X_k-\mu)\to0$ a.s.
\begin{solution}
不妨设 $\mu=0$。仿照定理~\ref{thm:random-series:slln} 的证明。令 $Y_k=X_k\one_{\{\abs{X_k}\leqslant k\}}$。

\emph{截断。}由引理~\ref{lem:strong-laws:truncation} 的证明，只有有限个 $k$ 使 $X_k\neq Y_k$，所以 $\frac1n\sum_{k\leqslant n}c_k(X_k-Y_k)\to0$ a.s.

\emph{随机部分。}$c_k(Y_k-EY_k)/k$ 独立，期望为零，$\sum_k\Var(c_kY_k/k)\leqslant C^2\sum_kEY_k^2/k^2\leqslant4C^2E\abs{X_1}$（引理~\ref{lem:strong-laws:variance}）。由定理~\ref{thm:random-series:convergence} 与 Kronecker 引理，$\frac1n\sum_{k\leqslant n}c_k(Y_k-EY_k)\to0$ a.s.

\emph{期望部分。}$\abs{c_kEY_k}\leqslant C\abs{EY_k}$，$EY_k\to\mu=0$（控制收敛），Cesàro 平均 $\frac1n\sum_{k\leqslant n}c_kEY_k\to0$。

三部分相加即得。注意 Etemadi 的证明在这里不好用：$\sum_{k\leqslant n}c_kX_k$ 关于 $n$ 不单调，“子列加单调夹逼”失效。随机级数方法不需要单调性。
\end{solution}
\end{exercises}

\begin{probleme}{B}{非负独立随机变量之和}
设 $X_n\geqslant0$ 独立。本题证明下列三条等价：
\[
\text{(i)}\ \sum_nX_n<\infty\ \text{a.s.};\quad\text{(ii)}\ \sum_n\bigl[P(X_n>1)+E(X_n\one_{\{X_n\leqslant1\}})\bigr]<\infty;\quad\text{(iii)}\ \sum_nE\frac{X_n}{1+X_n}<\infty .
\]
令 $Y_n=X_n\one_{\{X_n\leqslant1\}}$。
\begin{questions}
\item 证明 (ii) 与 (iii) 等价。
\begin{solution}
对 $x\in[0,1]$，$\frac x2\leqslant\frac x{1+x}\leqslant x$；对 $x>1$，$\frac12<\frac x{1+x}<1$。所以
\[
\frac12\bigl[P(X_n>1)+EY_n\bigr]\leqslant E\frac{X_n}{1+X_n}\leqslant P(X_n>1)+EY_n,
\]
两个级数同时收敛或同时发散。
\end{solution}
\item 证明 (ii) 蕴含 (i)。（不需要三级数定理。）
\begin{solution}
由 $\sum_nP(X_n>1)<\infty$ 与第一 Borel--Cantelli 引理，几乎必然只有有限个 $n$ 使 $X_n\neq Y_n$。由逐项积分，$E\sum_nY_n=\sum_nEY_n<\infty$，所以 $\sum_nY_n<\infty$ a.s.。于是 $\sum_nX_n$ 与 $\sum_nY_n$ 只差有限项，几乎必然有限。
\end{solution}
\item 设 (i) 成立。证明 $\sum_nP(X_n>1)<\infty$。
\begin{solution}
由 (i)，$X_n\to0$ a.s.，所以 $P(X_n>1\ \text{i.o.})=0$。事件 $\{X_n>1\}$ 独立，由第二 Borel--Cantelli 引理，$\sum_nP(X_n>1)<\infty$。
\end{solution}
\item 设 (i) 成立。证明 $\sum_nEY_n<\infty$。（提示：考虑 $E\exp(-\sum_nY_n)$。）
\begin{solution}
由 (i)，$0\leqslant\sum_nY_n\leqslant\sum_nX_n<\infty$ a.s.，所以 $e^{-\sum_nY_n}>0$ a.s.，$E e^{-\sum_nY_n}>0$（非负随机变量期望为零则几乎处处为零）。另一方面，由独立性，$E e^{-\sum_{n\leqslant N}Y_n}=\prod_{n\leqslant N}Ee^{-Y_n}$；令 $N\to\infty$，左边由有界收敛定理趋于 $Ee^{-\sum_nY_n}$。

对 $y\in[0,1]$，由 $e^{-y}$ 的凸性，$e^{-y}\leqslant1-(1-e^{-1})y\leqslant\exp(-cy)$，$c=1-e^{-1}$。所以 $Ee^{-Y_n}\leqslant\exp(-cEY_n)$（先用第一个不等式取期望，再用 $1-t\leqslant e^{-t}$），
\[
0<Ee^{-\sum_nY_n}=\lim_N\prod_{n\leqslant N}Ee^{-Y_n}\leqslant\exp\Bigl(-c\sum_nEY_n\Bigr).
\]
若 $\sum_nEY_n=\infty$，右边为 $0$，矛盾。由 B3)、B4)，(ii) 成立。

这里用 Laplace 变换 $Ee^{-\sum Y_n}$ 代替了三级数定理的必要性部分（后者要用中心极限定理），这是处理非负随机变量的常用办法。
\end{solution}
\end{questions}
\end{probleme}

\thinkingpart

\begin{probleme}[问题]{C}{随机符号级数的必要性}
设 $\varepsilon_n$ 独立，各以概率 $\frac12$ 取 $\pm1$，$c_n$ 是实数。本题不用中心极限定理证明：若 $\sum_n\varepsilon_nc_n$ 几乎必然收敛，则 $\sum_nc_n^2<\infty$。
\begin{questions}
\item 证明：若 $\sum_n\varepsilon_nc_n$ 几乎必然收敛，则 $c_n\to0$，从而可以设 $\abs{c_n}\leqslant1$。
\begin{solution}
收敛级数的项趋于零，$\abs{\varepsilon_nc_n}=\abs{c_n}$，所以 $c_n\to0$。去掉有限项不影响收敛性，也不影响 $\sum c_n^2$ 是否有限，所以可以设 $\abs{c_n}\leqslant1$。
\end{solution}
\item 设 $S=\sum_n\varepsilon_nc_n$（几乎必然收敛），$S_N=\sum_{n\leqslant N}\varepsilon_nc_n$。证明对实数 $t$，$E\cos(tS)=\prod_{n\geqslant1}\cos(tc_n)$。
\begin{solution}
对两个独立的 $\pm$ 号对称的随机变量 $U,V$（$U\eqd-U$，$V\eqd-V$），$E\cos(t(U+V))=E\cos(tU)E\cos(tV)-E\sin(tU)E\sin(tV)=E\cos(tU)E\cos(tV)$（由和角公式与独立性，奇函数 $\sin$ 的期望为零）。$S_{N-1}$ 与 $\varepsilon_Nc_N$ 独立且分布对称，对 $N$ 归纳得 $E\cos(tS_N)=\prod_{n\leqslant N}\cos(tc_n)$。令 $N\to\infty$，由有界收敛定理，$E\cos(tS)=\prod_n\cos(tc_n)$（部分积的极限）。
\end{solution}
\item 由 $S$ 几乎必然有限证明 $E\cos(tS)\to1$（$t\to0$），从而对充分小的 $t>0$，$\prod_n\cos(tc_n)>0$。再由 $\cos x\leqslant\exp(-x^2/4)$（$\abs x\leqslant1$）证明 $\sum_nc_n^2<\infty$。
\begin{solution}
$\cos(tS)\to1$ 逐点（$S$ 有限），有界，由有界收敛定理 $E\cos(tS)\to1$。取 $t\in(0,1]$ 使 $E\cos(tS)>\frac12$。$\abs{tc_n}\leqslant1$，由 $\cos x\leqslant1-\frac{x^2}2+\frac{x^4}{24}\leqslant1-\frac{x^2}4\leqslant e^{-x^2/4}$（$\abs x\leqslant1$；第一步是 Taylor 展开的余项估计），
\[
\frac12<\prod_n\cos(tc_n)\leqslant\exp\Bigl(-\frac{t^2}4\sum_nc_n^2\Bigr),
\]
所以 $\sum_nc_n^2<\infty$。结合\egref{ex:random-series:random-signs}：$\sum_n\varepsilon_nc_n$ 几乎必然收敛，当且仅当 $\sum_nc_n^2<\infty$。特别地，$\alpha\leqslant\frac12$ 时 $\sum\pm n^{-\alpha}$ 几乎必然不收敛（由 $0$-$1$ 律，不收敛的概率是 $1$）。

这个证明中的 $E\cos(tS)$ 其实是 $S$ 的特征函数，以后特征函数一讲的主角。
\end{solution}
\end{questions}
\end{probleme}

\exercisepart{课外知识：随机级数的和的分布}

\begin{probleme}{D}{$\sum\pm2^{-n}$ 与 $\sum\pm3^{-n}$}
设 $\varepsilon_n$ 独立，各以概率 $\frac12$ 取 $\pm1$。
\begin{questions}
\item 证明 $U=\sum_{n\geqslant1}\varepsilon_n2^{-n}$ 服从 $[-1,1]$ 上的均匀分布。
\begin{solution}
令 $\xi_n=\frac{1+\varepsilon_n}2\in\{0,1\}$，它们是独立的公平硬币。$U=\sum_n(2\xi_n-1)2^{-n}=2\sum_n\xi_n2^{-n}-1$。由定理~\ref{thm:distributions:coins-to-uniform}，$\sum_n\xi_n2^{-n}$ 在 $[0,1]$ 上均匀分布，所以 $U$ 在 $[-1,1]$ 上均匀分布。
\end{solution}
\item 证明 $V=\sum_{n\geqslant1}\varepsilon_n3^{-n}$ 的分布是奇异连续的。
\begin{solution}
同理 $V=\sum_n(2\xi_n-1)3^{-n}=\sum_n2\xi_n3^{-n}-\frac12$。$\sum_n2\xi_n3^{-n}$ 服从 Cantor 分布（第~\ref{lec:distributions} 讲作业 B），所以 $V$ 的分布是 Cantor 分布平移 $-\frac12$，奇异连续（\egref{ex:ac:cantor-coins}）。
\end{solution}
\item 说明：$\sum_n\varepsilon_nc_n$ 的分布的类型（离散、绝对连续、奇异）不能只从 $\sum c_n^2<\infty$ 看出来。
\begin{solution}
两个级数都满足 $\sum c_n^2<\infty$，几乎必然收敛，可是一个和是均匀分布，另一个是奇异的。一般地，由第~\ref{lec:absolute-continuity} 讲的分解定理与 $0$-$1$ 律的想法（Jessen--Wintner 定理），这样的级数的和的分布一定是“纯”的类型：纯离散、纯绝对连续或纯奇异连续，但究竟是哪一种可能非常难判断。比如 $\sum\pm\lambda^n$（$\frac12<\lambda<1$）对哪些 $\lambda$ 绝对连续，是一个至今仍有研究的问题（Bernoulli 卷积）。
\end{solution}
\end{questions}
\end{probleme}
